\honest{Beautiful Probability}{Eliezer Yudkowsky}{2008}

Should we expect rationality to be simple at some deep level? I begin with a complaint of ``the late great Bayesian Master,'' E. T. Jaynes. Two medical researchers try the same treatment. One decides to stop after 100 patients. The other decides to go on until the data show a cure rate ``definitely greater than 60\%.'' Both stop with 70 cures in 100 patients (and, presumably, equal control groups). Should we draw different conclusions?

Old-fashioned statistics treats these as different experiments, since they could have ended with different data, and may call only the first ``statistically significant.'' \nb{The figures bear this out. At a true rate of 60\%, 70 or more cures in a fixed 100 patients has probability 0.025. If the second researcher tests at the 2.5\% level after each patient, the chance of stopping by patient 100 is about 0.11.}

I let the non-Bayesians shrug that different tools give different results. \nb{The frequentist reason is not about tools. A significance level describes a procedure. Run on a treatment that cures exactly 60\%, the second rule declares a higher rate 11\% of the time by patient 100 and 29\% by patient 10,000. Bayesians answer, with Savage, that sampling to an appreciably false conclusion is practically impossible: at a true rate of 50\%, the rule succeeds about 1\% of the time.}

The Bayesian reply: ``Excuse you?'' The evidence cannot depend on ``the researcher's private thoughts,'' and the frequentists accuse Bayesians of being too subjective? How likely each state of Nature makes the data has nothing to do with intentions. The likelihood of the data is the same under both designs, so the conclusion must be the same, and at least one of the Old Style methods ``must discard relevant information - or simply do the wrong calculation.'' \nb{That step assumes the likelihood principle, that all the evidence is in the likelihood. Frequentists reject it, and its best-known derivation, Birnbaum's, is still disputed.}

Bayesians expect probability theory to be math: consistent, as Jaynes requires, with every theorem compatible with every other however derived. Compute 10 + 10 any legal way and you get 20. If you get 19 once, you made an illegal step, usually a division by zero, or in probability an infinity not taken as a limit. A real contradiction would bring down the whole edifice, set theory with it. Probability theory is also unique: Cox's theorem shows that any representation of uncertainty meeting certain constraints must map onto probability. Whatever is not Bayesian must fail a coherence test and can be Dutch-booked, accepting combinations of bets that are sure losses or refusing sure gains. \nb{The Dutch-book argument concerns degrees of belief. A significance test gives its hypothesis none, and the post does not say what bets it would lose.}

But should rationality be math? The world is messy, exact Bayesian calculation is often intractable, so why not many tools? Frequentists think in tools, and Bayesians in laws. I quote no statistician who holds the view I describe. Looking for laws is not looking for pretty tools; the second law of thermodynamics is not a pretty refrigerator. No engine running between two heat reservoirs beats a Carnot engine, and no real car engine is one, as no tire is a perfect cylinder. It would be absurd to conclude that thermodynamics does not apply to cars.

Likewise, approximations work to the extent that they approximate the Bayesian calculation and fail to the extent they depart, as a 747 is atoms even if we cannot compute its aerodynamics atom by atom. Their successes and failures are explainable in Bayesian terms, even if no one knows the explanation. Least-squares regression is the best point estimate under a Gaussian likelihood and a flat prior. Regularized regression adds a Gaussian prior on the weights.

When the exact Bayesian calculation can be done, ``You will never find a statistical method that yields a better answer,'' only cheaper ones, unless the other method uses knowledge, perhaps a disguised prior, that you left out. Put that in, and Bayes is again as good or better. With an ad-hoc tool someone may invent a cleverer one tomorrow. With the Bayes-optimal calculation you are done, like fitting a Carnot engine into your car. \nb{That holds on average over the calculation's own prior. Frequentists judge a method at each possible true value, and there another rule can do better. The complete class theorems, which the post does not cite, support a weaker claim: a method that no other beats everywhere is a Bayes rule for some prior.}

The toolboxers, it seems to me, look at the first differences of the cubes, 7, 19, 37, 61, and the Bayesians at the third, 6, 6, 6. Stability lies below the surface. Needing approximations does not change the law. The approximation is in the map, not the territory. Knowing the second law helps an engineer get close to ideal efficiency. The beauty of Bayesian theorems is a side effect of their being laws. An addendum points to chapter 37 of David MacKay's book for a fuller treatment of the opening problem.
